\section{Validering av provider-arkitekturen}
\label{sec:provider-validering}

For å undersøke om provider-arkitekturen faktisk var utvidbar,
implementerte vi to tillegsprovidere etter at AntMedia- og
NHN-integrasjonene var ferdige. \texttt{MockVideoProvider} var en enkel
stub uten eksterne avhengigheter, mens \texttt{DailyProvider} var en
reell integrasjon mot Daily.co sitt REST-API. Ingen av dem inngikk i HDOs
kravspesifikasjon; de ble kun brukt til arkitekturvalidering.

Endringsomfanget ble målt med \texttt{git diff --shortstat} mot
\texttt{develop}, med merge-base \texttt{5dd99b5} som felles baseline
for alle tre branchene. Med kjerneendringer mener vi endringer i
\texttt{RoomsController}, \texttt{IVideoProvider} eller de offentlige
modellene \texttt{VideoRoom}, \texttt{CreateRoomRequest} og
\texttt{WebRtcLinkResponse}. Provider-klasser, konfigurasjon,
DI-registrering og tester ble ikke regnet som kjerneendringer, siden
slike filer alltid må legges til når en ny provider innføres.
\newpage
\begin{table}[H]
\centering
\begin{tabular}{lrrrr}
\toprule
\textbf{Branch} & \textbf{Filer} & \textbf{Linjer +} &
\textbf{Linjer -} & \textbf{Kjerneendringer} \\
\midrule
\texttt{feature/mock-provider}       &  6 &  249 & 2 & 0 \\
\texttt{feature/daily-provider}      & 12 &  779 & 7 & 0 \\
\texttt{feature/provider-validation} & 14 & 1024 & 8 & 0 \\
\bottomrule
\end{tabular}
\caption{Endringsomfang per branch målt mot \texttt{develop}
(merge-base \texttt{5dd99b5}). \texttt{feature/provider-validation}
samler begge tilleggsproviderne.}
\label{tab:provider-validering}
\end{table}

Tabell~\ref{tab:provider-validering} viser at ingen av de tre
branchene berørte controller-laget eller API-kontrakten.
\texttt{git diff --name-status} bekreftet at \texttt{RoomsController},
\texttt{IVideoProvider} og de offentlige modellene var uendret i alle
branchene. Tallene gjelder lokal branch-HEAD mot \texttt{develop} på
måletidspunktet. E2E-skriptet ble heller ikke kjørt mot tilleggsproviderne.

% Alle branchene ble bygget og testet med \texttt{dotnet build} og
% \texttt{dotnet test}. På valideringsgrenen
% \texttt{feature/provider-validation} passerte 166 tester, mot 133 på
% leveransegrenen; økningen skyldtes tester for de to
% tilleggsproviderne og tilhørende DI-/factory-registrering.

% E2E-skriptet ble ikke kjørt mot tilleggsproviderne. Målet var ikke å
% dokumentere fullstendige Daily.co- eller mock-integrasjoner, men å
% undersøke om kjernelogikken forble uendret ved innføring av en ny
% provider. Resultatet støtter arkitekturpåstanden: backend-spesifikk
% logikk kunne isoleres i provider-laget mens controller-laget og
% API-kontrakten forble uberørt.

Målet var ikke å dokumentere fullstendige Daily.co- eller
mock-integrasjoner, men å undersøke om kjernelogikken forble uendret
ved innføring av en ny provider. Det er verdt å merke seg at
\texttt{MockVideoProvider} er en triviell stub, og at \texttt{DailyProvider}
representerer én reell REST-basert integrasjon. Funnene er dermed et
indisium på at isolasjonen fungerer i praksis, men utvalget er for lite
til å trekke generelle konklusjoner om arkitekturens skalerbarhet på
tvers av alle mulige provider-typer. Resultatet støtter likevel
arkitekturpåstanden for de implementerte tilfellene: backend-spesifikk
logikk kunne isoleres i provider-laget mens controller-laget og
API-kontrakten forble uberørt.
